% GENERATED FILE. Edit canonical lessons, then run npm run generate:books.
% canonical-source-hash: fnv1a64:986cb01f3f903451
% canonical-lessons: ML-C152-eli, ML-C152-tavala, ML-C152-ama, ML-C152-mayil, ML-C152-tatta

\chapter{Rat, Frog, Tortoise, Peacock, Parrot}
\label{ch:ml-rat-frog-tortoise-peacock-parrot}
\hlchaptermodality{152}

\begin{chapteropening}
\textbf{By the end of this chapter:} \emph{I can say a rat, a frog, a tortoise, a peacock and a parrot.}

Complete the last lesson of chapter 152: i can say a rat, a frog, a tortoise, a peacock and a parrot.
\end{chapteropening}

\section[eli]{\ml{എലി} (eli) — a rat, a mouse}

\label{lesson:ML-C152-eli}



\begin{quote}
Before the new one: say the Malayalam for a pig, then the Malayalam for a lion.

Say \textquotedblleft{}good morning,\textquotedblright{} and name the neighbour that uses the same word. (\emph{Suprabhātaṁ}; Telugu.)
\end{quote}

\subsection*{You'll want to know: \ml{എലി}}
\textbf{\ml{എലി}} — \emph{eli} — \textquotedblleft{}a rat, a mouse\textquotedblright{}.

\begin{grammarlens}[title={What you've built}]
One more word for this chapter.
\end{grammarlens}

\subsection*{Your turn}
\begin{itemize}
\raggedright
  \item \emph{Say it:} \emph{eli}
  \item \emph{Say it:} \emph{eli}, once more
  \item \emph{Say it:} say \emph{eli}
  \item \emph{Recall:} say the Malayalam for a pig, then the Malayalam for a lion, then say \emph{eli} again
\end{itemize}

\begin{tcolorbox}[breakable,colback=teal!4,colframe=teal!35!black,title={Before you move on}]
What does \ml{എലി} mean? (\textquotedblleft{}A rat, a mouse\textquotedblright{}.) Say it once more.
\end{tcolorbox}
\section[tavaḷa]{\ml{തവള} (tavaḷa) — a frog}

\label{lesson:ML-C152-tavala}



\begin{quote}
Before the new one: say the Malayalam for a lion, then the Malayalam for a rat.
\end{quote}

\subsection*{You'll want to know: \ml{തവള}}
\textbf{\ml{തവള}} — \emph{tavaḷa} — \textquotedblleft{}a frog\textquotedblright{}.

\begin{grammarlens}[title={What you've built}]
One more word for this chapter.
\end{grammarlens}

\subsection*{Your turn}
\begin{itemize}
\raggedright
  \item \emph{Say it:} \emph{tavaḷa}
  \item \emph{Say it:} \emph{tavaḷa}, once more
  \item \emph{Say it:} say \emph{tavaḷa}
  \item \emph{Recall:} say the Malayalam for a lion, then the Malayalam for a rat, then say \emph{tavaḷa} again
\end{itemize}

\begin{tcolorbox}[breakable,colback=teal!4,colframe=teal!35!black,title={Before you move on}]
What does \ml{തവള} mean? (\textquotedblleft{}A frog\textquotedblright{}.) Say it once more.
\end{tcolorbox}
\section[āma]{\ml{ആമ} (āma) — a tortoise}

\label{lesson:ML-C152-ama}



\begin{quote}
Before the new one: say the Malayalam for a rat, then the Malayalam for a frog.
\end{quote}

\subsection*{You'll want to know: \ml{ആമ}}
\textbf{\ml{ആമ}} — \emph{āma} — \textquotedblleft{}a tortoise\textquotedblright{}.

\begin{grammarlens}[title={What you've built}]
One more word for this chapter.
\end{grammarlens}

\subsection*{Your turn}
\begin{itemize}
\raggedright
  \item \emph{Say it:} \emph{āma}
  \item \emph{Say it:} \emph{āma}, once more
  \item \emph{Say it:} say \emph{āma}
  \item \emph{Recall:} say the Malayalam for a rat, then the Malayalam for a frog, then say \emph{āma} again
\end{itemize}

\begin{tcolorbox}[breakable,colback=teal!4,colframe=teal!35!black,title={Before you move on}]
What does \ml{ആമ} mean? (\textquotedblleft{}A tortoise\textquotedblright{}.) Say it once more.
\end{tcolorbox}
\section[mayil]{\ml{മയിൽ} (mayil) — a peacock}

\label{lesson:ML-C152-mayil}



\begin{quote}
Before the new one: say the Malayalam for a frog, then the Malayalam for a tortoise.

What do the two Sanskrit parts of \emph{sāyāhnaṁ} mean? (\textquotedblleft{}End\textquotedblright{} and \textquotedblleft{}day.\textquotedblright{})
\end{quote}

\subsection*{You'll want to know: \ml{മയിൽ}}
\textbf{\ml{മയിൽ}} — \emph{mayil} — \textquotedblleft{}a peacock\textquotedblright{}.

\begin{grammarlens}[title={What you've built}]
One more word for this chapter.
\end{grammarlens}

\subsection*{Your turn}
\begin{itemize}
\raggedright
  \item \emph{Say it:} \emph{mayil}
  \item \emph{Say it:} \emph{mayil}, once more
  \item \emph{Say it:} say \emph{mayil}
  \item \emph{Recall:} say the Malayalam for a frog, then the Malayalam for a tortoise, then say \emph{mayil} again
\end{itemize}

\begin{tcolorbox}[breakable,colback=teal!4,colframe=teal!35!black,title={Before you move on}]
What does \ml{മയിൽ} mean? (\textquotedblleft{}A peacock\textquotedblright{}.) Say it once more.
\end{tcolorbox}
\section[tatta]{\ml{തത്ത} (tatta) — a parrot}

\label{lesson:ML-C152-tatta}



\begin{quote}
Before the new one: say the Malayalam for a tortoise, then the Malayalam for a peacock.
\end{quote}

\subsection*{You'll want to know: \ml{തത്ത}}
\textbf{\ml{തത്ത}} — \emph{tatta} — \textquotedblleft{}a parrot\textquotedblright{}.

\begin{grammarlens}[title={What you've built}]
One more word for this chapter.
\end{grammarlens}

\subsection*{Your turn}
\begin{itemize}
\raggedright
  \item \emph{Say it:} \emph{tatta}
  \item \emph{Say it:} \emph{tatta}, once more
  \item \emph{Say it:} say \emph{tatta}
  \item \emph{Recall:} say the Malayalam for a tortoise, then the Malayalam for a peacock, then say \emph{tatta} again
\end{itemize}

\begin{tcolorbox}[breakable,colback=teal!4,colframe=teal!35!black,title={Before you move on}]
What does \ml{തത്ത} mean? (\textquotedblleft{}A parrot\textquotedblright{}.) Say it once more.
\end{tcolorbox}
